\section{A mass-increasing trajectory}\label{sec:increase}

We prove Theorem~\ref{thm:introincrease} by computing a positive one-sided
limit of the exact mass derivative. All metric and phase conventions
remain those of Section~\ref{sec:coefficients}.
For the initial matrices in that theorem,
$D^\dagger D=\diag(1/25,4)$ and
$A^\dagger A=\left(\begin{smallmatrix}1&2\\2&4\end{smallmatrix}\right)$
do not commute.

\subsection{The real limiting coefficient}
Put
\[
 b_1=1/25,\quad b_2=4,\quad p=e^{-b_1},\quad q=e^{-b_2},
 \quad R_i=\frac{1-e^{-b_i}}{b_i},\quad f_i=R_i-e^{-b_i},
\]
and define
\[
 \ell=\frac{p-q}{b_2-b_1},\qquad
 s_*=\frac{\sqrt{b_1b_2}(R_1-R_2)}{b_2-b_1}.
\]
At $\epsilon=0$ the coefficient blocks have the real symmetric limits
\begin{equation}\label{eq:realblocks}
 H=\begin{pmatrix}2p&2\ell\\2\ell&5q\end{pmatrix},\quad
 G=\begin{pmatrix}f_1-p&2s_*\\2s_*&4f_2-q\end{pmatrix},\quad
 w=4-R_1-4R_2=25p+q-22.
\end{equation}
The elementary bounds below give
\begin{equation}\label{eq:realblocksigns}
 \begin{gathered}
 -0.05060<\det H<-0.05059,\qquad
 -0.85991<\det G<-0.85990,\\
 2.03805<w<2.03806 .
 \end{gathered}
\end{equation}
Thus $H$ and $G$ each have one negative and one positive eigenvalue.
For a real symmetric $2$ by $2$ matrix $M$ with this property,
\begin{equation}\label{eq:negativeprojection}
 P_-(M)=\frac12\left(\id-
 \frac{2M-(\tr M)\id}{\sqrt{(\tr M)^2-4\det M}}\right).
\end{equation}
Write $P_V=P_-(H)$ and $P_U=P_-(G)$.

All limiting blocks are invertible. The polar unitary is continuous
near them. On the degree-zero branch its eigenvalues approach $1$
and $-1$ through the upper semicircle; the corresponding phases approach
$0$ and $\pi$. On the degree-one block apply the same observation to
$-\Omega_U$ and subtract $\pi$. It follows that
\begin{equation}\label{eq:phaselimit}
 \Phi_V\longrightarrow\pi P_V,\qquad
 \Phi_U\longrightarrow-\pi P_U,\qquad
 \Phi_W\longrightarrow0
 \quad(\epsilon\downarrow0).
\end{equation}
The argument is continuous on the closed upper semicircle with these
endpoint values. This assertion does not use continuity across a branch
cut from the lower semicircle.

\subsection{A finite real formula for the derivative}
Divide the limiting map velocities by $\pi$ and denote this directional
derivative by a prime. Equation~\eqref{eq:mapflow} becomes
\begin{equation}\label{eq:realvelocity}
 D'=\tfrac12(DP_V+P_UD),\qquad A'=\tfrac12AP_V .
\end{equation}
Since the central-charge trace is constant, the trace-norm derivative
formula gives
\begin{equation}\label{eq:positivelimit}
 \lim_{\epsilon\downarrow0}\frac{\dot m_\epsilon(0)}{\pi}
 =-2\tr(P_VH'+P_UG').
\end{equation}
Indeed, at the real invertible blocks the polar unitaries are their
sign matrices $\id-2P_V$, $\id-2P_U$, and $1$.
The identity contributions sum to the derivative of \eqref{eq:central},
which vanishes.

Here is a scalar-arithmetic evaluation of the right side of
\eqref{eq:positivelimit}. It avoids differentiating eigenvectors of
the coefficient. Let $B=\diag(b_1,b_2)$ and
$C=\left(\begin{smallmatrix}1&2\\2&4\end{smallmatrix}\right)$.
Set
\begin{equation}\label{eq:Bprime}
 B'=\tfrac12(P_VB+BP_V)+DP_UD,\quad
 \gamma=\frac{B'_{12}}{b_2-b_1},\quad
 K=\begin{pmatrix}0&\gamma\\-\gamma&0\end{pmatrix}.
\end{equation}
Use instantaneous singular-vector frames for $D$.
In the source frame, $B'+[B,K]$ is diagonal and
$b_i'=B'_{ii}$. The arrow Gram derivative is
\begin{equation}\label{eq:Cprime}
 \widetilde C'=\tfrac12(P_VC+CP_V)+[C,K].
\end{equation}
The corresponding target frame makes $D$ diagonal to first order.
Frame rotations do not change the traces
$\tr(P_VH')$ and $\tr(P_UG')$, since the projections commute with
the respective limiting blocks.

The remaining derivatives are
\begin{align}
 R_i'&=-\frac{f_i}{b_i}b_i',\qquad
 f_i'=\left(-\frac{f_i}{b_i}+e^{-b_i}\right)b_i',\label{eq:Rprime}\\
 \ell'&=\frac{(-pb_1'+qb_2')(b_2-b_1)
                    -(p-q)(b_2'-b_1')}{(b_2-b_1)^2},\label{eq:ellprime}\\
 s_*'&=\frac{s_*}{2}\left(\frac{b_1'}{b_1}+\frac{b_2'}{b_2}\right)\notag\\
 &\quad+\sqrt{b_1b_2}\,
  \frac{(R_1'-R_2')(b_2-b_1)-(R_1-R_2)(b_2'-b_1')}
       {(b_2-b_1)^2}.\label{eq:sprime}
\end{align}
In these frames,
\begin{align}
 H'_{ii}&=e^{-b_i}\bigl(\widetilde C'_{ii}-(1+C_{ii})b_i'\bigr),
 &H'_{12}&=\ell\widetilde C'_{12}+C_{12}\ell',\label{eq:Hprime}\\
 G'_{ii}&=f_i\widetilde C'_{ii}
 +\left[C_{ii}\left(-f_i/b_i+e^{-b_i}\right)+e^{-b_i}\right]b_i',
 &&\label{eq:Gprime}\\
 G'_{12}&=s_*\widetilde C'_{12}+C_{12}s_*'.&&\notag
\end{align}
Equations \eqref{eq:realblocks}--\eqref{eq:Gprime} are a finite expression
using rational operations, two exponentials and positive square roots.

\subsection{Strict positivity}
The two exponentials lie in the rational intervals
\begin{equation}\label{eq:expintervals}
 \begin{split}
 0.96078943915232320&<p<0.96078943915232322,\\
 0.01831563888873417&<q<0.01831563888873419.
 \end{split}
\end{equation}
To verify them, use the Taylor sums of degrees $31$ and $30$ for
$e^{-1/25}$, and of degrees $81$ and $80$ for $e^{-4}$.
The odd and even sums bound the exponential from below and above,
respectively, by the sign of the Taylor remainder. Each comparison with
the displayed decimal endpoint is a rational comparison.

Substitution in the finite formulas above, with interval operations,
gives the enclosures in Table~\ref{tab:intervals}. Each printed matrix
entry differs from its exact value by less than $10^{-8}$.
These deliberately coarse enclosures suffice to give
\begin{equation}\label{eq:coarsepositive}
 0.2395<-2\tr(P_VH'+P_UG')<0.2397.
\end{equation}
One may verify the table using only rational endpoints: every positive
square root can be enclosed by rational lower and upper bounds whose
squares enclose its argument. Thus the sign calculation does not
depend on a floating-point eigensolver or a finite-difference derivative.\footnote{The
source distribution includes two independent interval computations:
\path{figures/code/compute05_interval.py} and
\path{figures/code/compute05_independent.py}. The latter uses only
integer interval arithmetic from the Python standard library. They
enclose the one-sided limiting derivative, not a specified positive
value of $\epsilon$.}

\begin{table}[ht]
\caption{Enclosures for the sign calculation. For each symmetric matrix,
the entries are listed in the order $(11,12,22)$; each entry has error
less than $10^{-8}$.}\label{tab:intervals}
\centering
\begin{adjustbox}{max width=\textwidth,center}
\begin{tabular}{lrrr}
\toprule
Matrix & $11$ & $12$ & $22$\\
\midrule
$H$   & $1.92157888$ & $0.47599687$ & $0.09157819$\\
$G$   & $-0.94131486$ & $0.14845312$ & $0.89010617$\\
$P_V$ & $0.05643063$ & $-0.23075142$ & $0.94356937$\\
$P_U$ & $0.99355619$ & $-0.08001432$ & $0.00644381$\\
$H'$  & $0.01353176$ & $-0.24643619$ & $-0.29654199$\\
$G'$  & $0.06216737$ & $0.10644365$ & $0.12306918$\\
\bottomrule
\end{tabular}
\end{adjustbox}
\end{table}

\begin{proof}[Proof of Theorem~\ref{thm:introincrease}]
The coefficients, their directional derivatives, and the polar
unitaries extend continuously to the invertible real limiting blocks.
Equation~\eqref{eq:phaselimit} gives convergence of the actual flow
directions from $\epsilon>0$. The trace-norm derivative is consequently
continuous in this one-sided limit. Equations~\eqref{eq:positivelimit}
and \eqref{eq:coarsepositive} show that the limit is positive.
It follows that $\dot m_\epsilon(0)>0$ for all sufficiently small
$\epsilon>0$, when both central charges are admissible.
\end{proof}

The example disproves universal mass dissipation for the specified
positive phase, rather than merely exposing a failure of normality.
It does not preclude eventual cancellation, which
Theorem~\ref{thm:introcancellation} proves for these same complexes.
No change of vector field is used in that theorem.
